% GENERATED FILE. Edit canonical lessons, then run npm run generate:books.
% canonical-source-hash: fnv1a64:bc7b72aec5293f0a
% canonical-lessons: LA-C536-blandus, LA-C536-blanditia, LA-C536-suavis, LA-C536-deliciae, LA-C536-donum, LA-R536-review-529-536

\chapter{Charm and sweet words}
\label{ch:la-charm-and-sweet-words}
\hlchaptermodality{536}

\begin{chapteropening}
\textbf{By the end of this chapter:} \emph{I can talk about charm, sweetness and gifts.}

Complete the last lesson of chapter 536: I can talk about charm, sweetness and gifts.
\end{chapteropening}

\section[blandus, blanda, blandum]{blandus, blanda, blandum — coaxing, flattering; charming}

\label{lesson:LA-C536-blandus}



\begin{quote}
Before the new one: say the Latin for an admirer, then the Latin for a suitor.
\end{quote}

\subsection*{You'll want to know: blandus, blanda, blandum}
\textbf{blandus, blanda, blandum} — \textquotedblleft{}coaxing, flattering; charming\textquotedblright{}.

\begin{quote}
Fēlēs blanda ad dominam venit.
\end{quote}

\begin{quote}
\emph{The coaxing cat comes to her mistress.}
\end{quote}

Soft and winning words, true or false. English \emph{bland}, whose meaning has drifted.

\begin{cousinweb}
\textbf{blandus} is a native Latin adjective; Lewis \& Short suggests a link with \emph{mollis} \textquotedblleft{}soft\textquotedblright{} and English \emph{mild}, but the origin is not settled. English \emph{bland} comes from it.
\end{cousinweb}

\begin{grammarlens}[title={What you've built}]
One more word for telling, explaining and describing at B1.
\end{grammarlens}

\subsection*{Your turn}
\begin{itemize}
\raggedright
  \item \emph{Say it:} \emph{blandus, blanda, blandum}
  \item \emph{Say it:} \emph{blandus, blanda, blandum}, once more
  \item \emph{Say it:} a sentence of your own with \emph{blandus, blanda, blandum}
  \item \emph{Recall:} say the Latin for an admirer, then the Latin for a suitor, then say \emph{blandus, blanda, blandum} again
\end{itemize}

\begin{tcolorbox}[breakable,colback=teal!4,colframe=teal!35!black,title={Before you move on}]
What does blandus, blanda, blandum mean? (\textquotedblleft{}Coaxing, flattering; charming\textquotedblright{}.) Say it once more.
\end{tcolorbox}
\section[blanditia, blanditiae]{blanditia, blanditiae — flattery, coaxing}

\label{lesson:LA-C536-blanditia}



\begin{quote}
Before the new one: say the Latin for a suitor, then the Latin for coaxing.
\end{quote}

\subsection*{You'll want to know: blanditia, blanditiae}
\textbf{blanditia, blanditiae} — \textquotedblleft{}flattery, coaxing\textquotedblright{}.

\begin{quote}
Blanditiīs fīliae pater nōn resistit.
\end{quote}

\begin{quote}
\emph{The father does not resist his daughter's coaxing.}
\end{quote}

Feminine, first declension; often plural \emph{blanditiae} — coaxing words.

\begin{cousinweb}
\textbf{blanditia} is built on \emph{blandus} \textquotedblleft{}coaxing, flattering\textquotedblright{} (Lewis \& Short). English \emph{blandishment} comes from the related verb \emph{blandior}.
\end{cousinweb}

\begin{grammarlens}[title={What you've built}]
One more word for telling, explaining and describing at B1.
\end{grammarlens}

\subsection*{Your turn}
\begin{itemize}
\raggedright
  \item \emph{Say it:} \emph{blanditia, blanditiae}
  \item \emph{Say it:} \emph{blanditia, blanditiae}, once more
  \item \emph{Say it:} a sentence of your own with \emph{blanditia, blanditiae}
  \item \emph{Recall:} say the Latin for a suitor, then the Latin for coaxing, then say \emph{blanditia, blanditiae} again
\end{itemize}

\begin{tcolorbox}[breakable,colback=teal!4,colframe=teal!35!black,title={Before you move on}]
What does blanditia, blanditiae mean? (\textquotedblleft{}Flattery, coaxing\textquotedblright{}.) Say it once more.
\end{tcolorbox}
\section[suāvis, suāve]{suāvis, suāve — sweet, pleasant}

\label{lesson:LA-C536-suavis}



\begin{quote}
Before the new one: say the Latin for coaxing, then the Latin for flattery.
\end{quote}

\subsection*{You'll want to know: suāvis, suāve}
\textbf{suāvis, suāve} — \textquotedblleft{}sweet, pleasant\textquotedblright{}.

\begin{quote}
Odor rosārum suāvis est.
\end{quote}

\begin{quote}
\emph{The scent of the roses is sweet.}
\end{quote}

Two endings. Sweet to every sense, and in manner. English \emph{suave}.

\begin{cousinweb}
\textbf{suāvis} is inherited through Proto-Italic from a Proto-Indo-European word for \textquotedblleft{}sweet\textquotedblright{} (EtymDB), and English \emph{sweet} is a cousin. English \emph{suave} comes from it.
\end{cousinweb}

\begin{grammarlens}[title={What you've built}]
One more word for telling, explaining and describing at B1.
\end{grammarlens}

\subsection*{Your turn}
\begin{itemize}
\raggedright
  \item \emph{Say it:} \emph{suāvis, suāve}
  \item \emph{Say it:} \emph{suāvis, suāve}, once more
  \item \emph{Say it:} a sentence of your own with \emph{suāvis, suāve}
  \item \emph{Recall:} say the Latin for coaxing, then the Latin for flattery, then say \emph{suāvis, suāve} again
\end{itemize}

\begin{tcolorbox}[breakable,colback=teal!4,colframe=teal!35!black,title={Before you move on}]
What does suāvis, suāve mean? (\textquotedblleft{}Sweet, pleasant\textquotedblright{}.) Say it once more.
\end{tcolorbox}
\section[dēliciae, dēliciārum]{dēliciae, dēliciārum — a darling, pet; delight}

\label{lesson:LA-C536-deliciae}



\begin{quote}
Before the new one: say the Latin for flattery, then the Latin for sweet.
\end{quote}

\subsection*{You'll want to know: dēliciae, dēliciārum}
\textbf{dēliciae, dēliciārum} — \textquotedblleft{}a darling, pet; delight\textquotedblright{}.

\begin{quote}
Hic catulus dēliciae puellae est.
\end{quote}

\begin{quote}
\emph{This puppy is the girl's darling.}
\end{quote}

Feminine plural, first declension, often with a singular meaning: \emph{meae dēliciae!} — my darling! English \emph{delicious}.

\begin{cousinweb}
\textbf{dēliciae} is built on \emph{dēliciō} \textquotedblleft{}to allure\textquotedblright{}, from \emph{dē-} and \emph{laciō} \textquotedblleft{}to entice\textquotedblright{} (EtymDB; Lewis \& Short). English \emph{delicious} and \emph{delicacy} come from the family.
\end{cousinweb}

\begin{grammarlens}[title={What you've built}]
One more word for telling, explaining and describing at B1.
\end{grammarlens}

\subsection*{Your turn}
\begin{itemize}
\raggedright
  \item \emph{Say it:} \emph{dēliciae, dēliciārum}
  \item \emph{Say it:} \emph{dēliciae, dēliciārum}, once more
  \item \emph{Say it:} a sentence of your own with \emph{dēliciae, dēliciārum}
  \item \emph{Recall:} say the Latin for flattery, then the Latin for sweet, then say \emph{dēliciae, dēliciārum} again
\end{itemize}

\begin{tcolorbox}[breakable,colback=teal!4,colframe=teal!35!black,title={Before you move on}]
What does dēliciae, dēliciārum mean? (\textquotedblleft{}A darling, pet; delight\textquotedblright{}.) Say it once more.
\end{tcolorbox}
\section[dōnum, dōnī]{dōnum, dōnī — a gift}

\label{lesson:LA-C536-donum}



\begin{quote}
Before the new one: say the Latin for sweet, then the Latin for a darling.
\end{quote}

\subsection*{You'll want to know: dōnum, dōnī}
\textbf{dōnum, dōnī} — \textquotedblleft{}a gift\textquotedblright{}.

\begin{quote}
Amīcus mihi dōnum pulchrum mittit.
\end{quote}

\begin{quote}
\emph{My friend sends me a beautiful gift.}
\end{quote}

Neuter, second declension, from \emph{dō} — to give. English \emph{donation}.

\begin{cousinweb}
\textbf{dōnum} is inherited from a Proto-Indo-European word for \textquotedblleft{}gift\textquotedblright{}, built on the root of \emph{dō} \textquotedblleft{}to give\textquotedblright{} (EtymDB). English \emph{donor} and \emph{donation} come from the family.
\end{cousinweb}

\begin{grammarlens}[title={What you've built}]
One more word for telling, explaining and describing at B1.
\end{grammarlens}

\subsection*{Your turn}
\begin{itemize}
\raggedright
  \item \emph{Say it:} \emph{dōnum, dōnī}
  \item \emph{Say it:} \emph{dōnum, dōnī}, once more
  \item \emph{Say it:} a sentence of your own with \emph{dōnum, dōnī}
  \item \emph{Recall:} say the Latin for sweet, then the Latin for a darling, then say \emph{dōnum, dōnī} again
\end{itemize}

\begin{tcolorbox}[breakable,colback=teal!4,colframe=teal!35!black,title={Before you move on}]
What does dōnum, dōnī mean? (\textquotedblleft{}A gift\textquotedblright{}.) Say it once more.
\end{tcolorbox}
\section[(dialogue)]{Words from chapters 529–536 — a review}

\label{lesson:LA-R536-review-529-536}



\begin{quote}
Say the Latin for a stepmother and a stepfather.
\end{quote}

\subsection*{Your turn}
Say these aloud:

\begin{itemize}
\raggedright
  \item a stepmother, a stepfather, a stepson, a little daughter, brotherly
  \item a master, a home-born slave, a young maidservant, of a slave, a footman
  \item a steward, a doorkeeper, a porter, a litter-bearer, a farm housekeeper
  \item the freeing of a slave, a freedman, the rod of freedom, free-born, a felt cap
  \item a close friend, close friendship, a close bond, to hold dear, lovable
  \item a kiss, a loving kiss, to kiss, a kiss on the lips, an embrace
  \item to embrace, to clasp, a lover, an admirer, a suitor
  \item coaxing, flattery, sweet, a darling, a gift
\end{itemize}

\begin{tcolorbox}[breakable,colback=teal!4,colframe=teal!35!black,title={Before you move on}]
A stepmother? (\textbf{noverca, novercae}.) A stepfather? (\textbf{vitricus, vitricī}.) A stepson? (\textbf{prīvignus, prīvignī}.) A little daughter? (\textbf{fīliola, fīliolae}.) Brotherly? (\textbf{frāternus, frāterna, frāternum}.) A master? (\textbf{erus, erī}.) A home-born slave? (\textbf{verna, vernae}.) A young maidservant? (\textbf{ancillula, ancillulae}.) Of a slave? (\textbf{servīlis, servīle}.) A footman? (\textbf{pedisequus, pedisequī}.) A steward? (\textbf{ātriēnsis, ātriēnsis}.) A doorkeeper? (\textbf{iānitor, iānitōris}.) A porter? (\textbf{ostiārius, ostiāriī}.) A litter-bearer? (\textbf{lectīcārius, lectīcāriī}.) A farm housekeeper? (\textbf{vīlica, vīlicae}.) The freeing of a slave? (\textbf{manūmissiō, manūmissiōnis}.) A freedman? (\textbf{lībertīnus, lībertīnī}.) The rod of freedom? (\textbf{vindicta, vindictae}.) Free-born? (\textbf{ingenuus, ingenua, ingenuum}.) A felt cap? (\textbf{pilleus, pilleī}.) A close friend? (\textbf{familiāris, familiāris}.) Close friendship? (\textbf{familiāritās, familiāritātis}.) A close bond? (\textbf{necessitūdō, necessitūdinis}.) To hold dear? (\textbf{dīligō, dīligere}.) Lovable? (\textbf{amābilis, amābile}.) A kiss? (\textbf{ōsculum, ōsculī}.) A loving kiss? (\textbf{bāsium, bāsiī}.) To kiss? (\textbf{bāsiō, bāsiāre}.) A kiss on the lips? (\textbf{suāvium, suāviī}.) An embrace? (\textbf{amplexus, amplexūs}.) To embrace? (\textbf{amplector, amplectī}.) To clasp? (\textbf{complector, complectī}.) A lover? (\textbf{amāns, amantis}.) An admirer? (\textbf{amātor, amātōris}.) A suitor? (\textbf{procus, procī}.) Coaxing? (\textbf{blandus, blanda, blandum}.) Flattery? (\textbf{blanditia, blanditiae}.) Sweet? (\textbf{suāvis, suāve}.) A darling? (\textbf{dēliciae, dēliciārum}.) A gift? (\textbf{dōnum, dōnī}.)
\end{tcolorbox}
